\section{GitOps \& CI/CD}
{
    \metroset{titleformat frame=smallcaps}

% --- Slide 1 (words) ---
\begin{frame}{From Manual Deployment to Automated GitOps}
\begin{itemize}
    \item \textbf{The old way:} custom Dockerfiles built by hand, \texttt{docker-compose} for local runs, and manifests applied manually with \texttt{kubectl}.
    \item \textbf{The problem:} error-prone, not repeatable, and impossible to audit.
    \item \textbf{The monorepo:} one repository holds the code, Dockerfiles and Kubernetes manifests together.
    \item \textcolor{HUAblue}{\textbf{The goal:}} a single \texttt{git push} should take code all the way to the cluster, automatically.
\end{itemize}
\end{frame}

% --- Slide 2 (words) ---
\begin{frame}{The Pipeline: Jenkins (CI) and ArgoCD (CD)}
\begin{itemize}
    \item \textbf{Continuous Integration -- Jenkins:} on every push, a \texttt{Jenkinsfile} runs tests, builds the Docker images and pushes them to Docker Hub.
    \item \textbf{Continuous Delivery -- ArgoCD:} ArgoCD watches the Git repository and automatically syncs the cluster to match it.
    \item \textcolor{HUAblue}{\textbf{GitOps principle:}} Git is the single source of truth -- whatever is in Git is what runs in the cluster.
\end{itemize}
\end{frame}

% --- Slide 3 (diagram) ---
\begin{frame}{The GitOps Pipeline, End to End}
    \centering
    %\includegraphics[width=\textwidth]{figures/cicd-pipeline}
\end{frame}

}